%! Linear Functions: Slope, Forms \& Graphing — Unit 1, Lesson 1.1 — Lesson Plan
% Gradual-release plan (user direction, 2026-08-31): warm-up → guided notes & practice →
% group activity → debrief → close. No tiers, no exit ticket; homework is generated in-packet.
\documentclass[10pt]{article}
\usepackage{algebra2-article}
\usepackage{algebra2-boxes}
\usepackage{graphicx}   % -article does NOT load graphicx; needed for thumbnails/screenshots
\graphicspath{{images/}}
\newcommand{\TallMath}[1]{$\displaystyle #1\rule[-1.4em]{0pt}{3.2em}$}
\newcommand{\CourseName}{Algebra 2}
\newcommand{\MeetingLength}{60 minutes}
\newcommand{\UnitNumberName}{Unit 1: Linear Functions \quad}
\newcommand{\LessonNumberName}{Lesson 1.1: Linear Functions: Slope, Forms \& Graphing}

\begin{document}
\begin{center}
    {\Large\bfseries \CourseName \\
    {\small\bfseries \UnitNumberName \LessonNumberName}}
\end{center}

\begin{tcolorbox}[colback=forestbg, colframe=forest, boxrule=0.9pt, arc=2mm,
    left=2mm, right=2mm, top=1.5mm, bottom=1.5mm]
    \textbf{Primary Objective:} Students can compute \emph{slope} as a rate of change from a graph
    or from any two points, \emph{write} the equation of a line from a graph, from two points, or
    from a slope and a point, and move fluently among the three algebraic forms ---
    \emph{slope-intercept} $y=mx+b$, \emph{point-slope} $y-y_1=m(x-x_1)$, and \emph{standard}
    $Ax+By=C$ --- recognising that all three describe the \emph{same line}. They read a line as a
    transformation of the parent $y=x$ (a stretch or reflection changes the slope; a vertical shift
    changes the $y$-intercept), and identify parallel, perpendicular, horizontal, and vertical
    cases. Lesson~1.0 \emph{read} these features off a graph; today students \emph{produce} the
    equation.
    \\[2pt]
    \textbf{Standards (2023 VA SOL):} the Algebra~1 linear cluster \textbf{A.F.1}, reactivated as
    the foundation of Unit~1: \textbf{A.F.1a} (slope and intercepts as rate of change),
    \textbf{A.F.1b} (transformations of the parent $y=x$ change slope and $y$-intercept),
    \textbf{A.F.1c} (write and interpret slope-intercept, point-slope, and standard forms),
    \textbf{A.F.1d} (write the equation of a line from a graph, from two points, or from a slope
    and a point), and \textbf{A.F.1e} (parallel and perpendicular lines). It launches the
    transformational lens that \textbf{A2.F.1} extends to every later function family.
    \\[2pt]
    \textbf{Lesson model:} \emph{Gradual release} --- I do, we do, you do together, you do alone. A
    warm-up activates prior knowledge; \textbf{Guided Notes} build the three forms with the class on
    one worked context (a phone plan); a \textbf{Group Activity} applies them to two climbing gyms,
    a comparison, and a model; a \textbf{debrief} consolidates; \textbf{homework} is the independent
    practice.
\end{tcolorbox}

\renewcommand{\arraystretch}{1.4}
\begin{skillbox}[Priority Ideas \& Skills]{goldbox}
\begin{tabularx}{\linewidth}{@{}X|X@{}}
\textbf{Priority Ideas \& Skills} & \textbf{Key Understandings (the ``why'')} \\[2pt]
\hline
\vspace{-6pt}
\begin{itemize}[leftmargin=1.1em, nosep, after=\vspace{2pt}]
  \item Find the \textbf{slope} from a graph or from any two points.
  \item Write a line in \textbf{slope-intercept} form from a graph or from a rate and a start.
  \item Write a line in \textbf{point-slope} form from a slope and \emph{any} point.
  \item Convert point-slope $\to$ slope-intercept $\to$ \textbf{standard}, and back.
  \item Describe $y=mx+b$ as a transformation of the parent $y=x$.
  \item Identify \textbf{parallel}, \textbf{perpendicular}, horizontal, and vertical lines.
\end{itemize}
&
\vspace{-6pt}
\begin{itemize}[leftmargin=1.1em, nosep, after=\vspace{2pt}]
  \item Any two points give the same slope, because a line's rate of change is \emph{constant} ---
        the defining property from Lesson~1.0.
  \item Point-slope exists because the $y$-intercept is often the one point you \emph{do not} have.
  \item \textbf{Target misconception:} that two equations which \emph{look} different must be
        different lines. Priya's $D-24=2(n-3)$ and the group's $D=2n+18$ are the same rule ---
        the reason the three forms are taught together.
  \item Every non-vertical line is $y=x$ after a stretch/reflection (slope) and a vertical shift
        (intercept); $x=a$ is the one line no $y=mx+b$ can express.
\end{itemize}
\end{tabularx}
\end{skillbox}

\begin{skillbox}[{Vocabulary, Concepts, \& Theorems}]{sky}
\renewcommand{\arraystretch}{1.3}
\begin{tabularx}{\linewidth}{@{}l X@{}}
\textbf{Slope / rate of change} & \TallMath{m=\frac{y_2-y_1}{x_2-x_1}=\frac{\Delta y}{\Delta x}}: the constant change in output per unit change in input. \\
\textbf{Slope-intercept form} & $y=mx+b$: $m$ is the slope, $(0,b)$ the $y$-intercept. Reads off a rate and a starting value. \\
\textbf{Point-slope form} & $y-y_1=m(x-x_1)$: built directly from a slope and \emph{any} known point $(x_1,y_1)$. \\
\textbf{Standard form} & $Ax+By=C$ with $A,B,C$ integers and $A\ge0$; convenient for intercepts. \\
\textbf{Parent function} & The simplest line $y=x$ (slope $1$, through the origin); every non-vertical line is a transformation of it. \\
\textbf{Transformation} & $kf(x)$ stretches/reflects (changes the slope); $f(x)+k$ shifts vertically (changes the $y$-intercept). \\
\textbf{Parallel / perpendicular} & Equal slopes / opposite-reciprocal slopes ($m_1m_2=-1$). \\
\end{tabularx}
\end{skillbox}

\begin{fixedskillbox}[Lesson at a Glance --- \MeetingLength]{forestbg}
\renewcommand{\arraystretch}{1.25}
\begin{tabularx}{\linewidth}{@{}l c X X@{}}
\textbf{Phase} & \textbf{Min} & \textbf{Students} & \textbf{Teacher} \\
\hline
Warm-Up & 5 & Three prior skills, alone then check with a neighbour & Circulate; debrief items~2--3 aloud (``two equations that looked different'') \\
Guided Notes & 15 & Fill the vocabulary and the four sections with the class & \emph{I do / we do}: build all three forms on \emph{one} context, the phone plan; four \texttt{work} blocks at the board \\
Group Activity & 25 & \emph{Two Receipts} parts 1--4 in groups of 3--4 & Circulate; questions, cues, prompts; choose work to display \\
Debrief & 10 & Share findings; correct their own pages & Put Priya's rule beside a group's; ``same line, different outfit'' \\
Close \& Homework & 5 & Record the assignment; preview 1.2 & Assign the homework page; name what changed today \\
\end{tabularx}
\end{fixedskillbox}

\begin{fixedskillbox}[Warm-Up --- Activate Prior Knowledge \& Spiral Review]{forestbg}
    The Warm-Up rehearses exactly what today leans on, without naming anything. \textbf{(1)} Counting
    steps between two marked points on a graphed line --- ``3 right, 4 up, so $\tfrac43$ up per
    step'' --- is the seed of \emph{slope from two points}; the follow-up (``where does it cross the
    vertical axis?'') is deliberately the awkward $-\tfrac73$, so the picture alone will not give
    it --- which is the argument for point-slope form. \textbf{(2)} Solving $2x+3y=12$ and $5x-y=8$
    for $y$ seeds \emph{standard $\to$ slope-intercept}. \textbf{(3)} Distributing $y-5=3(x-2)$ and
    $y+1=-2(x+4)$ and tidying seeds \emph{point-slope $\to$ slope-intercept} --- the exact algebra
    the notes and the activity both need. \textbf{Debrief items 2 and 3 aloud, briefly}, and land
    only this: \emph{two equations that looked nothing alike turned into something much easier to
    read.}
\end{fixedskillbox}

\begin{skillbox}[Hook]{forestbg}
    ``Your phone bill does not list the monthly fee --- only the total. Two months: \textbf{4~GB cost
    \$52}, \textbf{10~GB cost \$70}.'' Ask: \textbf{``How much more data, for how much more
    money?''} ($6$~GB for \$18.) \textbf{``So what does one GB cost?''} (\$3.) \textbf{``Then what is
    the flat monthly fee?''} (\$40 --- back off $4\times\$3$ from \$52.) Land the idea before naming
    anything: \emph{nobody handed you the starting value; you worked back to it.} Today we build the
    equation of a line from whatever we are given --- a graph, two points, or a slope and any point
    at all --- and we meet the form that lets you write it \emph{before} you know the intercept.
\end{skillbox}

\begin{skillbox}[Guided Notes --- the Lesson (15 min)]{forestbg}
\begin{multicols}{2}
\textbf{1. Slope from any two points.} $m=\frac{\Delta y}{\Delta x}=\frac{y_2-y_1}{x_2-x_1}$;
subtract in the same order top and bottom. Work the first \texttt{work} block on the hook's own
points $(4,52)$ and $(10,70)$: $m=3$. Ask why \emph{any} two points work --- answer from Lesson~1.0,
the rate is constant.

\textbf{2. Three ways to write one line.} Fill the three-row table. Then the sentence that matters:
\emph{point-slope exists because the $y$-intercept is often the number nobody gives you.} Write
$y-52=3(x-4)$ straight from the point you already have --- no intercept required.

\columnbreak
\textbf{3. Same line, different outfit.} Two \texttt{work} blocks, back to back: expand
$y-52=3(x-4)$ to $y=3x+40$ --- \emph{and there is the \$40 fee from the hook, falling out of the
algebra} --- then move $x$ across to $3x-y=-40$. Close by testing $x=10$ in both: \$70 either way.

\textbf{4. The parent, and four special cases.} $y=x$ with $y=2x-4$ on one grid: multiplying changes
the slope, adding shifts the intercept. Then the four-cell table --- parallel, perpendicular,
horizontal, vertical --- ending on the one line $y=mx+b$ cannot write.

\textbf{Practice (the ``we do'').} Slope $-\tfrac12$ through $(2,-5)$: point-slope, then a
\texttt{work} block to slope-intercept, then parallel and perpendicular slopes.
\end{multicols}
\end{skillbox}

\begin{skillbox}[Group Work --- \emph{Two Receipts} (25 min, groups of 3--4)]{redbox}
One common task, no tiers. \textbf{Launch (2 min):} ``Two climbing gyms charge the same way ---
something to walk in the door, then a price per climb. One shows you a graph. The other hands you two
receipts.'' One packet each, one shared conversation: \emph{for every number you write down, be ready
to show me where it came from.}
\begin{multicols}{2}
\textbf{What students do.} \textbf{Part~1} (Summit) is the easy case: read the \$6 walk-in fee off
the axis, get \$3 per climb from $(0,6)$ to $(4,18)$, assemble $C(n)=3n+6$, check at $n=7$.
\textbf{Part~2} (Basecamp) removes the graph \emph{and} the starting value: two receipts give \$2 per
climb, and groups work \emph{backwards} to the \$18 fee. Then the crux, \textbf{2c}: Priya wrote
$D-24=2(n-3)$ straight off her receipt --- name the form, test $n=7$, expand it. 2d asks for standard
form. \textbf{Part~3} compares the two gyms and finds where the answer switches ($n=12$, \$42), and
asks whether they are parallel. \textbf{Part~4} is the model, the drone descent, running
two-points $\to$ point-slope $\to$ slope-intercept $\to$ standard in one problem.

\columnbreak
\textbf{What the teacher does.} Circulate. Listen before speaking. Give a \emph{question, cue, or
prompt}, not an answer:
\begin{itemize}[leftmargin=1.1em, nosep]
  \item 1b: ``Finger on $(0,6)$, finger on $(4,18)$. How far up, over how many climbs? Now one climb.''
  \item 2b: ``Three climbs \emph{and} the fee cost \$24. What do the three climbs cost alone?''
  \item 2c: \emph{Do not settle it.} ``Try $n=7$ in both. Now $n=10$. What would convince you?''
  \item 3b: ``Which gym wins at $n=2$? At $n=15$? So somewhere in between, what happened?''
  \item 4c: ``Which number is the rate? Which is the start? So which one moved?''
\end{itemize}
Pick work to display: a group's \emph{backwards} work for the \$18, and --- essential --- a group
that has expanded Priya's rule.
\end{multicols}
\end{skillbox}

\boxguard
\begin{skillbox}[Debrief (10 min --- what goes on the board, in a second color)]{forestbg}
Do not re-teach the activity; land four things on the displayed student work.
\begin{enumerate}[leftmargin=1.3em, nosep, itemsep=2pt]
  \item \textbf{2b, the backwards move.} Put a group's arithmetic up: three climbs cost \$6, so the
        fee is \$18. That is the same reasoning as the phone plan's \$40 --- name the connection.
  \item \textbf{2c is the must-land moment.} Display Priya's rule beside a group's $D=2n+18$ and
        expand hers \emph{line by line} on the board. Say it out loud: \emph{same line, different
        outfit}. Then name why point-slope exists --- the intercept is the point you did not have.
  \item \textbf{Part~3 --- what the slopes tell you.} Different slopes, so not parallel, so they
        cross exactly once; $n=12$ is that crossing, and it is why the cheaper gym switches.
  \item \textbf{4c --- rate vs.\ start.} Starting from $60$~m changes $b$, not $m$: the line shifts
        up, it does not tilt. A student who changes $m$ is still fusing the two roles --- send them
        to the graph, not to the algebra.
\end{enumerate}
If time is short, cut item~1 and protect items~2 and~3.
\end{skillbox}

\boxguard
\begin{skillbox}[Active Monitoring --- Watch For]{redbox}
\begin{itemize}[leftmargin=1.2em, nosep]
  \item computing slope as $\frac{x_2-x_1}{y_2-y_1}$, or subtracting in a different order top and
        bottom (Warm-Up~1, notes~1, homework~1);
  \item in 2b, dividing \$24 by 3 to ``get the price per climb'' --- that folds the walk-in fee into
        the rate; push them back to the \emph{difference} between the two receipts;
  \item \textbf{2c: declaring Priya wrong because her equation looks different} --- the lesson's
        target misconception. Make them expand it, not vote on it;
  \item reading $y-y_1=m(x-x_1)$ as needing the $y$-intercept for $(x_1,y_1)$ (homework~2b);
  \item giving $x=4$ a slope of $0$ and $y=4$ an undefined one --- exactly backwards (homework~4b);
  \item calling $y=\tfrac12x-2$ perpendicular to $y=\tfrac12x-4$ because ``the numbers changed''
        (homework~6);
  \item in 2d and 4b, leaving standard form with a negative $A$ or a fraction.
\end{itemize}
Cold-call prompts: ``Where did that number come from --- which two points?'' ``Same line or a
different line? Convince me.'' ``Which number is the rate and which is the start?''
\end{skillbox}

\boxguard
\begin{skillbox}[Reinforcement \& Extension]{goldbox}
\textbf{Homework} (in the packet, collected next class): slope from two points three times, the last
two being the horizontal ($0$) and vertical (undefined) cases; one graphed line written
slope-intercept \emph{and} point-slope with a \texttt{work} block showing they agree --- Priya again,
on paper; converting $y+2=-\tfrac34(x-4)$ through slope-intercept to standard, fractions and all; the
transformation $y=x \to y=-3x+2$ plus the $y=4$ / $x=4$ boundary; a candle model; and an SOL-style
multiple-choice item on perpendicular slopes. \textbf{Extension:} a parallel-line construction from
standard form, and why point-slope gives the same line from \emph{any} point. \textbf{DeltaMath
override:} well covered there --- swap in a DeltaMath set on slope, forms of a line, and
parallel/perpendicular if you prefer auto-graded practice. \textbf{Preview:} Lesson~1.2,
\emph{Absolute Value Equations \& Inequalities} --- every line today crossed zero exactly once; next
comes a rule that reaches the same value \emph{twice}, because it measures distance.
\end{skillbox}

% ── Teacher notes — one per component, in packet order ──
\begin{teachernote}[Warm-Up]
Item~1's ``where does it cross the vertical axis?'' is deliberately $-\tfrac73$ --- ugly, and not
readable off the picture. Do not resolve it; let a group grumble, because that grumble is the
argument for point-slope form in notes section~2. Debrief only items 2 and 3, and only the
observation that two different-looking equations tidied into something readable. Five minutes; if it
runs long, drop the second equation in each of items 2 and 3.
\end{teachernote}

\begin{teachernote}[Guided Notes]
Fifteen minutes, paced by section: 3 for the vocabulary and hook, 3 for section~1, 3 for section~2,
4 for section~3, 2 for section~4. \textbf{Keep the whole notes block on one context} --- the phone
plan --- so students see the three forms as three outfits on \emph{one} line rather than three
unrelated procedures. The payoff is scripted: the \$40 the class worked back to in the hook falls out
of the section~3 \texttt{work} block as $b$. Four \texttt{work} blocks are printed here (slope; two
conversions; the practice box), and students write in the reserved space while the key prints the
same steps there, so the two files paginate alike. If you fall behind, cut section~4's transformation
graph to one sentence --- but do \emph{not} cut the four special cases; homework~4 needs them.
\end{teachernote}

\begin{teachernote}[Group Activity]
Budget roughly 5 minutes for Part~1, 9 for Part~2, 5 for Part~3, and 6 for Part~4. Part~1 is the easy
case and groups will move fast --- protect the time for Part~2, where the real work is. If a group
finishes Part~2 early, ask them to write Basecamp's rule starting from the \emph{other} receipt
($D-32=2(n-7)$) and check that it also expands to $2n+18$; that is the strongest version of the
argument. Common slips: dividing \$24 by 3 in 2a; writing standard form as $-2n+D=18$ (needs
$A\ge0$); and in 3c, saying the lines are parallel because ``they both go up.''
\end{teachernote}

\begin{teachernote}[Homework]
Item~1's last two entries separate students who have the horizontal/vertical cases straight: $y$ the
same both times gives $m=0$; $x$ the same both times gives a zero denominator, so \emph{undefined}.
Item~3 is the fraction-heavy conversion --- expect $y+2=-\tfrac34x+3$ to be mis-distributed as
$-\tfrac34x-3$. Item~5(b) asks for the burn-out time, $t=9.5$~h, which requires solving
$-2t+19=0$ --- not just reading the rule. Item~6 is the formative check: sort into ``chose B (got it)
/ chose A (used the same slope --- parallel, not perpendicular) / chose C (flipped the sign only) /
chose D (took the reciprocal only)'' to decide whether Lesson~1.2 opens with a one-minute re-look at
opposite reciprocals. Three \texttt{work} blocks (items 2, 3, and 5a) plus one in the extension.
\end{teachernote}

\end{document}
